% Generated draft: requires mathematical and source review.
\documentclass[11pt]{article}
\usepackage[margin=25mm]{geometry}
\usepackage{fontspec}
\setmainfont{Latin Modern Roman}
\usepackage{amsmath,amssymb,mathtools,unicode-math}
\setmathfont{Latin Modern Math}
\usepackage{enumitem,needspace,xurl,hyperref}
\setlength{\parindent}{0pt}
\setlength{\parskip}{5pt}
\setlength{\emergencystretch}{3em}
\newcommand{\N}{\mathbb{N}}
\newcommand{\Z}{\mathbb{Z}}
\newcommand{\Q}{\mathbb{Q}}
\newcommand{\R}{\mathbb{R}}
\newcommand{\C}{\mathbb{C}}
\newcommand{\F}{\mathbb{F}}
\newcommand{\A}{\mathbb{A}}
\newcommand{\PP}{\mathbb{P}}
\newcommand{\E}{\mathbb{E}}
\newcommand{\eps}{\varepsilon}
\newcommand{\dd}{\,\mathrm{d}}
\begin{document}
\section*{UnsolvedMath ID: 11000041}
\textbf{Source problem code:} catalogue entry
\subsection{Definitions and mathematical statement}
For \(g\geq3\), let \(\mathcal M_g\) be the moduli \emph{orbifold} of closed genus-\(g\) Riemann surfaces, modeled by \(\operatorname{Teich}_g/\operatorname{Mod}_g\) with its stabilizer data. An orbifold covering respects local uniformizing charts and isotropy; a finite covering has finitely many sheets.

The printed maximality conjecture says that \(\mathcal M_g\) does not finitely orbifold-cover any other smooth orbifold. In the usual no-proper-quotient interpretation, this asks that every connected finite orbifold quotient \(p:\mathcal M_g\to\mathcal O\) be trivial. Whether other means merely a nonisomorphic target, or excludes every covering of degree greater than \(1\), needs clarification.

The source also proposes a stronger statement: if \(N\to\mathcal M_g\) is any finite orbifold cover and \(N\to\mathcal O\) is an orbifold cover, then \(\mathcal O\) should itself orbifold-cover \(\mathcal M_g\). The second covering is not explicitly stated finite, and the text does not require that these maps form a compatible factorization of the given map \(N\to\mathcal M_g\). This stronger statement is explicitly presented only as something that may be true.

The smooth-category formulation has a material unresolved issue. Complex conjugation gives a natural involution on the moduli orbifold, or equivalently one may enlarge to the extended mapping class group; allowing that quotient produces a natural smooth orbifold quotient of degree \(2\). Thus the intended allowed category must explain whether such quotients are excluded, for example by additional complex/holomorphic restrictions. No such restriction is printed, and it is not silently inserted. The following related topology question is mentioned but not included.
\subsection{Short English statement}
Is the moduli orbifold maximal among the allowed finite orbifold quotients for \(g\geq3\)? The proposed stronger cover statement and the missing restriction excluding conjugation quotients require clarification.
\subsection{Sources}
\begin{itemize}
\item \emph{Problems on Mapping Class Groups and Related Topics}, edited by Farb, Conjecture 4.2, ``\(\mathcal M_g\) is maximal,'' supplied PDF page 38. \url{https://www.math.uchicago.edu/~farb/papers/mcgbook.pdf}. Recorded access 2026-07-29; extraction pdf-font-structure. No separate reference is supplied for this item.
\item Local UnsolvedMath registry, \texttt{data/unsolvedmath/problems.json}, permanent ID \(11000041\), exact source code \texttt{catalogue entry}; supplied OPEN-TRIAGE summary says status cannot be verified without the statement.
\end{itemize}
\end{document}
